\section{Introduction}

Ce document présente l'organisation pédagogique retenue pour la préparation pratique du
\textbf{Plongeur Niveau 1 (N1) --- Plongeur encadré à 20 m (PE20)}, conformément au manuel de
formation technique FFESSM (version décembre 2025).

\subsection{Rappel des prérogatives}

Le plongeur N1 est capable de réaliser des plongées d'exploration jusqu'à 20 m de profondeur,
au sein d'une palanquée, avec un guide de palanquée (GP) qui prend en charge la conduite de la
plongée. Ces plongées sont réalisées dans le cadre d'une organisation sécurisée, mise en place
par un directeur de plongée (DP), selon les règles définies par le Code du Sport.

\subsection{Organisation}

\begin{itemize}[leftmargin=1.2cm]
  \item \textbf{Durée totale} : environ 3 mois.
  \item \textbf{Rythme} : 1 séance par semaine.
  \item \textbf{Durée d'une séance} : 1 h à 1 h 30 selon le contenu.
  \item \textbf{Lieu} : piscine / fosse de plongée, dans l'espace de 0 à 6 m.
  \item \textbf{Encadrement} : enseignement et validation des compétences par un encadrant
        licencié au minimum E1. Conformément à l'article A.\,322-83 du Code du Sport, un
        plongeur en cours de formation technique peut évoluer dans l'espace de 0 à 20 m sous
        la responsabilité d'un encadrant E2. L'accoutumance à la profondeur doit être
        progressive.
  \item \textbf{Effectif} : formation en petit groupe / palanquée, afin de développer les
        notions d'entraide et de cohésion dès les premières séances.
\end{itemize}

\subsection{Modalités d'évaluation}

L'évaluation des compétences se fait en \textbf{contrôle continu}. Le niveau est validé
lorsque toutes les compétences du référentiel sont acquises. Les connaissances théoriques sont
évaluées à l'oral, lors des mises en situation pratiques : il n'y a pas d'examen écrit.
L'accent est mis sur la prévention.

Au CST, une \textbf{séance théorique} est prévue en \textbf{décembre}, complétée par un
\textbf{examen écrit}. Cet examen n'a pas vocation à valider le niveau : il permet
d'identifier les lacunes de chaque élève afin de cibler les révisions nécessaires pour
compléter sa formation.

La \textbf{fiche de suivi et d'évaluation} détaillée par sous-compétences (section
\ref{sec:fiche-evaluation}) permet à l'encadrant de suivre la progression de chaque élève
séance après séance.
